\documentclass[../../../include/open-logic-section]{subfiles}

\begin{document}
	
\olfileid{sth}{z}{nat}
\olsection{د خپلو طبيعي عددونو ټاکنه}

%Let us take it for granted that the informal principle \stagesinf{}
%can be justified. It is also, though, worth commenting on the
%particular axiom we extracted from that informal principle. 

په \olref[infinity-again]{defnomega} کښې مو د «طبيعي عددونو»
عبارت په ښکاره تعريف کړے دے. دا ټاکنه څنګه درک کول پکار دي؟
دا يوه مابعدالطبيعي ادعا نۀ ده، بلکې يوازې دا پرېکړه ده چې ځينې
سټونه د طبيعي عددونو په توګه \emph{واخلو}. د دې نظر پر وجوهاتو مو
په \olref[sfr][rel][ref]{sec}، \olref[sfr][arith][ref]{sec} او
\olref[sfr][infinite][dedekindsproof]{sec} کښې لږ بحث کړے و.
خو دا وجوهات لا هم دقيقولے شو.

زموږ د نامتناهيت بديهي اصل د \citet{VonNeumann1925} پيروي کوي.
خو دلته يو بل اصل دے چې پر ځاے ئې منلے شو:

\begin{defish}
\emph{د زرمېلو د \citeyear{Zermelo1908Untersuchungen} د نامتناهيت
بديهي اصل.} يو سټ $A$ شته چې $\emptyset \in A$ او
$(\forall x \in A)\{x\} \in A$ وي.
\end{defish}

که د خپل (له فون نويمان څخه الهام اخيستي) نامتناهيت اصل پر ځاے
مو د زرمېلو اصل کارولے و، نو په همدې ډول به د ډېډېکېنډ له مخې
يو نامتناهي سټ او ورسره د ډېډېکېنډ الجبر راکړل شوے و. د زرمېلو
په لاره کښې به زموږ د الجبر ځانګړے غړے بيا هم $\emptyset$ و
(زموږ د $0$ استازے)، خو يو پر يو تابع به د نگاشت
$x \mapsto \{x\}$ په وسيله ورکړل شوې وه، نۀ د
$x \mapsto x \cup \{x\}$ په وسيله. د دې تر ټولو ساده پايله دا ده
چې زرمېلو $2$ د $\{\{\emptyset\}\}$ په توګه اخلي، خو موږ
(له فون نويمان سره) $2$ د $\{\emptyset, \{\emptyset\}\}$ په توګه اخلو.

ولې د نامتناهيت يو اصل د بل پر ځاے غوره کړو؟ مهم عملي سبب دا
دے چې د فون نويمان لاره د نامتناهيت هاخوا عددونو د سمبالولو
دپاره ښه غځېدلے شي. دا به له \olref[ordinals][]{chap} څخه
وروسته وڅېړو. خو د \emph{حساب کولو} له ساده نظره، دواړه لارې
يو شان ګټورې دي. نو که څوک درته وايي چې طبيعي عددونه سټونه
\emph{دي}، څرګنده پوښتنه دا ده: \emph{کوم سټونه دي؟}

همدا دقيقه پوښتنه \citet{Benacerraf1965} مشهوره کړه.
خو دا ويل پکار دي چې دا د هغې ښکارندې يوازې تر ټولو مشهوره
بېلګه ده چې مخکې هم ډېر ځله ورسره مخامخ شوي يو. بنسټيزه خبره
دا ده: د سټ تيوري د ډېرو «شهودي» ډوله شيانو د \emph{تمثيل}
يوه لاره راکوي: حقيقي، ناطق، صحيح او طبيعي عددونه، هو؛ خو
مرتبه جوړې، تابعې او اړيکې هم. مګر د سټ تيوري د تمثيل يو
\emph{يوازينے} انتخاب هېڅکله نۀ راکوي. \emph{تل} داسې بديلونه
شته چې په نېغه به همدومره ښه کار راکړي.

\end{document}
